% Source: https://artofproblemsolving.com/community/c3h2452011_ash_ketchum_problems
% Stable identity: aops-topic-2452011-q22
Let $r=69$, let $O$ be the centre and $V_0,\dots,V_5$ the vertices in order. Each triangle $OV_kV_{k+1}$ is equilateral with side $r$.

\textbf{One leaf.} The circle with centre $V_0$ and radius $r$ passes through $O$ and $V_1$ (both at distance $r$ from $V_0$), and its arc from $V_1$ to $O$ subtends $\angle V_1V_0O=60^\circ$. The region between this arc and the chord $OV_1$ is a circular segment of area
\[\text{sector}-\text{triangle}=\dfrac{60}{360}\pi r^2-\dfrac{\sqrt3}{4}r^2=r^2\left(\dfrac{\pi}{6}-\dfrac{\sqrt3}{4}\right).\]
It bulges away from $V_0$, into triangle $OV_1V_2$. Likewise the arc of the circle centred at $V_2$ from $O$ to $V_1$ cuts off a congruent segment on the other side of $OV_1$. These two segments form the leaf from $O$ to $V_1$:

\begin{asy}
usepackage("amsmath");
size(5cm);
defaultpen(fontsize(10pt));
real r=1;
pair[] V; for(int k=0;k<6;++k) V.push(r*dir(60k));
pair O=(0,0);
for(int i=0;i<6;++i){
  path pet=arc(V[(i+5)%6],r,60i+60,60i+120)--arc(V[(i+1)%6],r,60i+240,60i+300)--cycle;
  fill(pet,rgb(0.85,0.95,0.85));}
path seg1=arc(V[0],r,120,180)--cycle;
path seg2=arc(V[2],r,300,360)--cycle;
fill(seg1,rgb(0.5,0.8,0.5)); fill(seg2,rgb(0.6,0.75,0.95));
for(int k=0;k<6;++k) draw(arc(V[k],r,60k+120,60k+240),darkgreen);
draw(V[0]--V[1]--V[2]--V[3]--V[4]--V[5]--cycle,linewidth(0.9));
draw(O--V[1],red); draw(V[0]--O,dashed); draw(V[0]--V[1],dashed); draw(V[2]--O,dashed);
dot(O); label("$O$",O,dir(255));
label("$V_1$",V[1],N); label("$V_0$",V[0],E); label("$V_2$",V[2],N);
label("$60^\circ$",V[0]+0.28*dir(150),fontsize(7pt));
\end{asy}

so each leaf has area $2r^2\left(\dfrac{\pi}{6}-\dfrac{\sqrt3}{4}\right)=r^2\left(\dfrac{\pi}{3}-\dfrac{\sqrt3}{2}\right)$.

\textbf{No overlaps.} Inside triangle $OV_1V_2$ lie two segments: the one on chord $OV_1$ (arc centred at $V_0$) and, by symmetry, the one on chord $OV_2$ (arc centred at $V_3$). The axis of symmetry of this triangle is the line through $O$ perpendicular to $V_0O$, so it is tangent at $O$ to the circle centred at $V_0$; hence the segment on $OV_1$ lies on $V_1$'s side of the axis, touching it only at $O$, and by symmetry the segment on $OV_2$ lies on $V_2$'s side. So the segments meet only at $O$. Hence the six leaves do not overlap, and together they cover $6r^2\left(\dfrac{\pi}{3}-\dfrac{\sqrt3}{2}\right)=r^2\left(2\pi-3\sqrt3\right)$.

\textbf{Uncovered area.} The hexagon has area $6\cdot\dfrac{\sqrt3}{4}r^2=\dfrac{3\sqrt3}{2}r^2$, so the uncovered area is
\[\dfrac{3\sqrt3}{2}r^2-\left(2\pi-3\sqrt3\right)r^2=\left(\dfrac{9\sqrt3}{2}-2\pi\right)r^2.\]
With $r^2=4761$ this is
\[\boxed{\dfrac{42849\sqrt3}{2}-9522\pi}\approx7194.1.\]
